% Pista ana-23s-q2 · Anàlisi
% Procedència: PAU Catalunya 2023 · Convocatòria de setembre · Sèrie 2 · Qüestió 2
\documentclass[11pt,a4paper]{article}
\usepackage{pista}
\pistainfo{Catalunya 2023}{Convocatòria de setembre}{Sèrie 2}

\begin{document}

\section*{\textcolor{blauPista}{Qüestió 2}}

\noindent Sigui la funció $f(x) = \dfrac{1}{x}$.

\subsection*{\textit{a)} \normalfont Calculeu l'equació de la recta tangent a la gràfica de la funció $f$ en el punt d'abscissa $x = 2$. \hspace*{\fill}\textnormal{[0,75~punts]}}

\pista{El pendent de la recta tangent que busques és $f'(2)$, i el punt és $\bigl(2, f(2)\bigr)$. Per tant:
\[
y - f(2) = f'(2)\,(x - 2)
\]
En primer lloc, calcula $f'(x)$. Recorda que $\frac{1}{x} = x^{-1}$: la pots derivar com una potència. A continuació, calcula $f(2)$ i $f'(2)$, i substitueix-los a la fórmula.}

\subsection*{\textit{b)} \normalfont Calculeu l'equació de la recta tangent a la gràfica de la funció $f$ en el punt d'abscissa $x = k$, en què $k$ és un nombre real positiu. \hspace*{\fill}\textnormal{[0,75~punts]}}

\pista{Fes el mateix que a l'apartat \textit{a)}, però amb $k$ en lloc de $2$:
\[
y - f(k) = f'(k)\,(x - k)
\]
Calcula $f(k)$ i $f'(k)$: l'equació de la tangent et quedarà en funció de $k$. Et convé simplificar-la i aïllar la $y$, perquè a l'apartat \textit{c)} la faràs servir.}

\subsection*{\textit{c)} \normalfont Comproveu que, tal com es pot veure en la figura de sota, la recta de l'apartat \textit{b} determina un triangle d'àrea constant amb els semieixos positius de coordenades. Calculeu aquesta àrea. \hspace*{\fill}\textnormal{[1~punt]}}

\pista{El triangle té els vèrtexs a l'origen i als dos punts on la recta tangent de l'apartat \textit{b)} talla els eixos. En primer lloc, troba aquests dos punts. Per al tall amb l'eix d'abscisses, fes $y = 0$ a l'equació de la recta tangent i resol l'equació; en canvi, per trobar el punt de tall amb l'eix d'ordenades és més senzill, perquè només has de fer $x = 0$.

\medskip
A continuació, el triangle és rectangle, i la base i l'altura són les distàncies de l'origen a aquests dos punts:
\[
\text{Àrea} = \frac{\text{base} \cdot \text{altura}}{2}
\]
Quan simplifiquis, la $k$ ha de desaparèixer. Al final, escriu explícitament que l'àrea no depèn de $k$: és el que et demanen comprovar.}

\end{document}
